\section{Introduction}
\label{sec:intro}
\IEEEPARstart{E}{arth} observation archives grow in volume, combine several
sensors and resolutions, and move to the cloud. Their use depends on a
storage layout with a global spatial reference, fast local access and
scalable chunking. Discrete Global Grid Systems (DGGS) divide the Earth into
hierarchically indexed cells that serve as a shared reference for storage,
aggregation and analysis~\cite{ogc_dggs,iso19170}; H3~\cite{h3},
S2~\cite{s2geometry} and HEALPix~\cite{gorski2005healpix} differ in cell
shape, refinement rule and area preservation. Cloud-native products move
in this direction: chunked array stores such as Zarr~\cite{moore2023zarr},
community tooling for DGGS data cubes such as
\texttt{xdggs}~\cite{kmoch2024xdggs}, and ESA's Earth Observation
Processing Framework, which republishes the Sentinel archive as Zarr.

HEALPix was designed for astronomical data on the
sphere~\cite{gorski2005healpix}. Its cells have equal area, its refinement
is dyadic, and in its nested index all descendants of a cell form one
contiguous block. These properties suit Earth observation as well, with
one practical limit: HEALPix resolution changes only by powers of two, so
the native sampling of an instrument almost never matches a level. The
\SI{10}{m} bands of Sentinel-2~\cite{drusch2012sentinel2} fall between
level 19 (\SI{12.4}{m}, too coarse) and level 20 (\SI{6.2}{m},
geometrically safe and nominally oversampled). This paper asks whether
that oversampling is expensive once storage and compression are both
counted. Level-20 samples are not independent observations at
\SI{6.2}{m}; they form an overcomplete representation whose information
content the sensor response still limits. If this redundancy compresses
well, storing the data at level 20 could cost no more than keeping the
same information on the native UTM grid. The native product has a cost
of its own: under an operational tiling layout such as the MGRS granules
of Sentinel-2 it carries padding, overlap and per-tile metadata beyond
the area needed, and this paper uses a $30\%$ overhead as an illustrative
scenario for one layout (Section~\ref{sec:overhead}).

The answer, in a data manager's terms, is given in
Section~\ref{sec:headline}. Before compression, a level-20 tile takes $2.58$
times the bytes of the area-matched native raster. After compression to
the same reconstruction error, it takes $2.06$ times as many when the
tolerated error is a quarter of a percent of a typical scene's contrast,
and $0.93$ times as many, a bare crossing, marginal in significance
(Section~\ref{sec:rd_main}), when it is a quarter of that contrast. This
holds for a tile filled by Clough--Tocher interpolation, which passes a
smooth surface through the observed pixels and reads it at the HEALPix
cell centres. The PSF-aware transfer of the companion
study~\cite{delouis2026resample}, which also undoes part of the instrument
blur, aims at a sharper latent-field estimand; its tile stays at
$2.3$--$2.6$ times the native raster at every error, bytes whose
reconstruction benefit is established by that study, not by this one. The
choice of transfer is therefore one of the largest sensitivities in the
storage cost of a HEALPix carrier.
Section~\ref{sec:glossary} defines the nine terms the paper relies on.

The question is answered empirically, with established compression
libraries and no HEALPix-specific coding. The finer HEALPix level is
treated as a plain 2-D raster, obtained by Morton (Z-order) de-interleaving
of the nested chunk, and passed to standard codecs as an ordinary image.
The comparison covers matched ground area: the centroid of the HEALPix
tile defines an axis-aligned, equal-area square cut from the native UTM
raster (Section~\ref{sec:utm_crop}), so the UTM side is always the
untouched native product. The contributions are:
\begin{itemize}
    \item a storage--distortion formulation of the resolution mismatch
    between native Sentinel-2 sampling and the next finer HEALPix level,
    with a geometric accounting: the $2.58\times$ oversampling of level 20
    makes the raw HEALPix tile the \emph{larger} representation by
    $\approx\!1.99$ under the illustrative $30\%$ overhead, so only
    compression can make HEALPix competitive;
    \item an equal-area rate--distortion protocol that separates the
    overhead-free carrier/native size ratio from the overhead of the native
    layout through a critical-overhead $\gamma_\star$ analysis, at matched
    error and matched budget, against a native UTM square of the same
    nominal area;
    \item an evaluation over 40 Sentinel-2 regions, four land-cover
    classes, latitudes from $43.5^\circ$S to $64.1^\circ$N (both HEALPix
    cell-shape regimes) and five general-purpose codecs, with sensitivity
    analyses on codec, target, transfer and assumed overhead, a
    multiple-testing correction that seven of the eight per-target
    comparisons survive, an analytic-field control showing that sampling
    geometry alone reproduces the decline of the gap with the target, and
    a smoothed-native control showing that transfer-induced smoothing is
    sufficient to account for a residual of the observed magnitude
    (Section~\ref{sec:support_controls});
    \item a reusable, assertion-guarded pipeline that shares its sites,
    products and resampling operator with the companion
    study~\cite{delouis2026resample}, so the storage and transfer questions
    are answered on the same data.
\end{itemize}

The competitive regime reported below belongs to the complete chain
and depends on four quantities: the overhead of the native product,
carried as $\gamma_\star(e)$ so that no value is assumed; the transfer
(Section~\ref{sec:transfer_cost}); the codec
(Section~\ref{sec:same_codec}); and the reconstruction reference
(Section~\ref{sec:common_reference}). These four conditions belong with
every statement of the form ``HEALPix is smaller''.
